% The document class supplies options to control rendering of some standard
% features in the result.  The goal is for uniform style, so some attention 
% to detail is *vital* with all fields.  Each field (i.e., text inside the
% curly braces below, so the MEng text inside {MEng} for instance) should 
% take into account the following:
%
% - author name       should be formatted as "FirstName LastName"
%   (not "Initial LastName" for example),
% - supervisor name   should be formatted as "Title FirstName LastName"
%   (where Title is "Dr." or "Prof." for example),
% - degree programme  should be "BSc", "MEng", "MSci", "MSc" or "PhD",
% - dissertation title should be correctly capitalised (plus you can have
%   an optional sub-title if appropriate, or leave this field blank),
% - dissertation type should be formatted as one of the following:
%   * for the MEng degree programme either "enterprise" or "research" to
%     reflect the stream,
%   * for the MSc  degree programme "$X/Y/Z$" for a project deemed to be
%     X%, Y% and Z% of type I, II and III.
% - year              should be formatted as a 4-digit year of submission
%   (so 2014 rather than the academic year, say 2013/14 say).

\documentclass[ % the name of the author
                    author={Rzan Alkabier},
                % the name of the supervisor
                    supervisor={Dr. Zahraa Abdallah},
                % the degree programme
                    degree={MSc},
                % the dissertation    title (which cannot be blank)
                     title={Denoising Extracellular Recording for Spike Detection and Classification on Single Channels},
                % the dissertation subtitle (which can    be blank)
                    subtitle={Does denoising improve detection and classification performance?},
                % the dissertation     type
                      type={},
                % the year of submission
                      year={2026}]{dissertation}

\begin{document}

% =============================================================================

% This section simply introduces the structural guidelines.  It can clearly
% be deleted (or commented out) if you use the file as a template for your
% own dissertation: everything following it is in the correct order to use 
% as is.
% =============================================================================

% This macro creates the standard UoB title page by using information drawn
% from the document class (meaning it is vital you select the correct degree 
% title and so on).


\maketitle

% After the title page (which is a special case in that it is not numbered)
% comes the front matter or preliminaries; this macro signals the start of
% such content, meaning the pages are numbered with Roman numerals.

\frontmatter

% This macro creates the standard UoB declaration; on the printed hard-copy,
% this must be physically signed by the author in the space indicated.

\makedecl

% LaTeX automatically generates a table of contents, plus associated lists 
% of figures, tables and algorithms.  The former is a compulsory part of the
% dissertation, but if you do not require the latter they can be suppressed
% by simply commenting out the associated macro.

\tableofcontents
\listoffigures
\listoftables
%\listofalgorithms
%\lstlistoflistings

% The following sections are part of the front matter, but are not generated
% automatically by LaTeX; the use of \chapter* means they are not numbered.

% -----------------------------------------------------------------------------
\chapter*{Ethics Statement}
\noindent
This project fits within the scope of the blanket ethics application 15208, as reviewed by my supervisor (Zahraa Abdallah)

\chapter*{Abstract}
Chronic pain is linked to the firing of C-fibre nociceptors, which can be recorded by microneurography. However, the microelectrode records the activity of multiple fibres at once, alongside other background noise. The spikes from a target fibre must first be detected and separated from those other spikes, which is challenging because of the similarity of C-fibre waveforms. Denoising is usually used on these low signal-to-noise-ratio recordings on the assumption that it improves spike detection and sorting. 
\bigskip

\noindent 
This thesis tests the hypothesis that denoising improves single-channel spike detection and classification in these recordings, and evaluates this under a protocol designed to avoid the information leakage that can inflate the results. A leakage-safe pipeline was developed, in which spike templates, denoisers, and detection parameters are all fitted on a development set and performance is scored once on a held-out evaluation dataset. Three denoisiers were used for this analysis [Wavelet, dense autoencoder and convolutional autoencoder], which were compared on five different channels independently, of a rat saphenous nerve recording on which the target fibre appears. 
\bigskip

\noindent
Contrary to the common assumption, denoising did not improve performance. All denoisers reduced noise but attenuated the spike, with the lowest denoiser being the most destructive; denoising did not reduce the false positives in detection; and both classifiers were at chance on every channel, before and after denoising. The consistency of results across channels suggests that single-channel target waveform shape can not be separated from co-active fibres and that denoising cannot overcome this limitation. This leakage-safe evaluation shows that denoising does not improve single-channel spike sorting and motivates multi-channel approaches that exploit spatial information.  


% -----------------------------------------------------------------------------
\chapter*{Supporting Technologies}
\begin{itemize}
    \item Python (version 3.13) was used as the main programming language for all data processing, model training, and analysis. 
    \item NumPy for numerical arrays, and SciPy for signal processing, specifically \texttt{scipy.signal.find\_peaks} for peak detection, \texttt{scipy.signal.correlate} for template matching, and \texttt{scipy.ndimage. uniform\_filter1d} for the local normalisation in the normalised cross-correlation. 
    \item PyTorch was used to implement and train the dense and convolutional denoising autoencoders. 
    \item PyWavelet was used for the stationary wavelet transform denoising
    \item scikit-learn was used for machine learning classifiers, specifically the one-class and two-class support vector machine \texttt{OneClassSVM}, \texttt{SVC}, the \texttt{StandardScaler}, and the evaluation metrics \texttt{roc\_auc\_score}, \texttt{average\_precision\_score}. 
    \item h5ph was used to read the MATLAB v7.3 (HDF5)files containing the data and the reference annotations. 
    \item Mtaplotlib was used for all figures and visualisations
    \item GitHub was used for version control (\url{https://github.com/Razan-88/Decoding_Chronic_Pain.git})
    \item \LaTeX\ was used to format this dissertation.  
    
\end{itemize}
% -----------------------------------------------------------------------------
\chapter*{Acknowledgements}
First and foremost, I would like to thank my supervisor, Dr Zahraa Abdallah. She provided excellent academic feedback, but more importantly, she showed incredible empathy when my personal life became difficult. Her patience and encouragement during my worst period of anxiety are the main reason I was able to complete this project. 

\noindent
I would also like to thank Omar Swehli for funding my studies and for his encouragement, and Dr Abdulrafiu for his endless support. 

\noindent
I'm deeply grateful to my family and friends for their constant support throughout this dissertation. In particular, I would like to thank my mum, Dr Jehan Alswaihli, for always being a role model and for inspiring me throughout my academic journey.

\noindent
Finally, I would like to thank my cat [Prince] for providing emotional support and keeping me company during the late hours of working. 
% =============================================================================

% After the front matter comes a number of chapters; under each chapter,
% sections, subsections and even subsubsections are permissible.  The
% pages in this part are numbered with Arabic numerals.  Note that:
%
% - A reference point can be marked using \label{XXX}, and then later
%   referred to via \ref{XXX}; for example Chapter\ref{chap:context}.
% - The chapters are presented here in one file; this can become hard
%   to manage.  An alternative is to save the content in seprate files
%   the use \input{XXX} to import it, which acts like the #include
%   directive in C.

\mainmatter

% -----------------------------------------------------------------------------

\chapter{Introduction: Contextual Background}
\label{chap:introduction}
Chronic pain is a common and complex condition that persists for longer than three months, affecting 30-33\% of the population worldwide and 1 in 3 adults in the UK \cite{shetty_chronic_2025}. A meta-analysis of 19 populations by \cite{fayaz_prevalence_2016} showed that 43.5\% of the UK population (28 million adults) experience chronic pain and 14.3\% (8 million adults) live with moderate or severe disabling chronic pain. Due to the ageing population in England, chronic pain is estimated to grow by 1.7 million by 2040 \cite{toby_watt_health_2023}. Patients with chronic pain can also experience impacts on their quality of life, daily living activities, and psychological well-being \cite{andersen_psychosocial_2014}. 
\bigskip

\noindent
The treatment for chronic pain is limited because the underlying cause of this condition is unknown. The standard treatments for acute pain are opioids, but they are not considered safe for long-term use for chronic pain due to the shifts in the physiology of the nociceptors. 
Chronic pain is caused by injury or disease in the nervous system and is associated with persistent hyperactivity of nociceptors, which are specialised to detect damaging stimuli and/or inflammation in rodents or humans \cite{walters_persistent_2023}. Therefore, a constant activation of these nociceptors, especially the C-fibre nociceptor, is a driver of prolonged pain and can lead to maladaptive neuroplasticity. Therefore, measuring the abnormal activities of these nociceptors could reveal and allow for precise pain targeting \cite{ossipov_central_2010}. 
\bigskip

\noindent
These nociceptor activities can be recorded using microneurography, a technique where an electrode is inserted through the skin into a nerve bundle to record action potentials from a single fibre of interest. However, the spike patterns of these C-fibre nociceptors in the microneurography recordings are challenging because of the small spike amplitude, low signal-to-noise ratio, the conduction velocity and activity dependence, and their waveform shapes are similar between different fibres \cite{troglio_hybrid_2026}. Comprehensive spike detection and sorting analysis are therefore required. 
\bigskip

\noindent
Many studies have researched and developed methods for extracellular spike detection and sorting for multi-electrode data, but not many focus on single-contacts where the spatial information is absent, and the discharge rates of unmyelinated peripheral nerve fibres are lower. Because the signal has a low signal-to-noise ratio, denoising is often applied as a preprocessing step, such as bandpass filtering, to improve detection and sorting of spikes. However, it is unclear if denoising actually improves spike detection and sorting; a denoiser may reduce noise, but distort the spike or introduce artefacts. For single-channel analysis, denoising is more difficult because the recording includes spikes from other fibres that denoisers can not remove.
\bigskip

\noindent
This thesis evaluates whether denoising improves single-channel spike detection and classification using chronological split, in which all tuning is performed on a development set, and performance is scored once on the held-out data. Three denoisers are compared: wavelet and two autoencoders, dense and convolutional autoencoders. The evaluation is repeated independently across five different channels that contain annotated spikes of interest to assess whether the findings are consistent across channels. 

\subsection{Aims and objectives}
\begin{itemize}
    \item to assess the quality of each denoiser, in terms of noise reduction and spike preservation relative to the raw signal; 
    \item to determine whether denoising improves spike detection;
    \item to determine whether denoising improves classification of the target fibre against non-target fibres; and 
    \item to assess whether these findings are consistent across channels. 
\end{itemize}

\subsection{Structure}
The following chapters are as follows: Chapter 2 explores the nature of extracellular recording and discusses different methodologies for spike detection and classification, signal denoising, the strengths and weaknesses of these methods for such recordings and the gap this work is trying to address. Chapter 3 introduces the studied dataset and the different channels being analysed; the spike detection and sorting pipeline executed and the different denoisers used for the signal, with key algorithms described briefly. Chapter 4 is the critical evaluation of spike detection and classification before and after denoising, with measuring the quality of different denoising methods and limitations of this analysis. Finally, Chapter 5 concludes the thesis and suggests future work. 


% -----------------------------------------------------------------------------

\chapter{Technical Background}
\label{chap:background}

\section{Extracellular Recordings} 
\label{sec:extracellular}

Extracellular recordings measure the voltage generated by the action potentials of a single or multiple neurons near an electrode placed in the extracellular space \cite{kechris_removing_2021, guo_extracellular_2020, lewicki_review_1998}. Extracellular action potentials have low-amplitude signals, usually lasting 1 to 2 ms with peak amplitude of $100\mu\text{V}$ or less, compared to that of intracellular \cite{hill_quality_2011}, which places them in proximity to the background noise floor. So, the signal is sampled at a high rate of 20 to 30 kHz.
\bigskip

\noindent
Because the electrode in the extracellular space does not only record a single fibre, but multiple fibre activities in its vicinity \cite{lewicki_review_1998}, separating these into the activity of individual units is the purpose of spike sorting (Section\ref{sec:detection}). In addition to the activities recorded from the target unit and other fibres, the recorded signal also contains noise from other sources such as thermal and electrical noise from the recording hardware, and low-amplitude activity of many distant units \cite{azami_extracellular_2015}. Because the recording contains a mixture of several fibre activities that are embedded in noise, distinguishing individual fibres, such as a target nociceptor from other activities, is not possible by reading the recording directly. The signal must be preprocessed to detect spikes and assign them to their fibre, which is addressed in the following sections.  

\section{Spike detection and sorting}
\label{sec:detection}
The activity of several sources is recorded, so spike sorting is essential to detect the action potentials and assign each to its unit before any firing pattern analysis \cite{rey_past_2015}. Therefore, spike sorting is divided into two stages: detection, which identifies when a spike occurs, and sorting or classification, which determines which unit produced it.  
\bigskip

\noindent
Most spike sorting analysis pipelines follow a common sequence \cite{rey_past_2015}: starting with band-pass filtering to isolate the spike frequency band; the spikes are then detected, usually by amplitude thresholding; then a waveform is extracted around each spike; and then features are computed for dimensionality reduction, for example by using principal component analysis (PCA) \cite{quiroga_unsupervised_2004}; and then the spikes are clustered in this feature space, where each cluster corresponds to a putative unit. According to \cite{troglio_hybrid_2026}, PCA reduces dimensionality but might discard discriminative details that raw waveforms retain, which was found to perform better. 
Overall, the feature-based spike sorting pipeline is based on the assumption that spikes from different units form distinctive clusters in the chosen feature space \cite{lewicki_review_1998}. This assumption could be true when the units have distinct waveform shapes, such as spikes from different spatial positions captured by multiple electrodes, yielding a distinct multi-channel signature \cite{harris_accuracy_2000}. Therefore, modern automated sorters like WaveClus \cite{quiroga_unsupervised_2004} and Kilosort and MountainSort produce sufficient separation. However, this assumption of cluster separation fails when waveforms are similar between units, such as those in single-channel recordings, which is discussed in the following section. 
\bigskip

\noindent
Spike detection is an important step, as the quality of the algorithm influences the performance of the next steps, such as spike sorting and classification, so any errors in detection of the numbers and location of the target spikes would inevitably propagate through all the later stages in the analysis \cite{azami_extracellular_2015}. Multiple reasons can make spike detection difficult; the recorded spikes are corrupted by non-target neuron activity, background noise and recording hardware noise. Moreover, the activities of nearby neurons in the recording may merge as noise that correlates with the target signal \cite{azami_extracellular_2015}.  

A single contact lacks this spatial information, so algorithms rely mostly on temporal voltage dynamics such as waveform amplitude, shape, latency, and frequency content. This involves setting a fixed voltage threshold, classifying any crossings as candidate spikes and extracting a short segment of the recorded voltage around these events for subsequent spike detection, which then can be extracted for waveform analysis, classification or spike sorting \cite{laboy-juarez_normalized_2019}. For a single spike not to be counted more than once as a separate event, a refractory period or a shadow period is enforced between the detected spike and the next spike \cite{laboy-juarez_normalized_2019}. 
\bigskip

\noindent
A fixed-amplitude threshold detection is simple and commonly used, but its performance depends on the user's selected threshold, so a threshold that is too high can cause low-amplitude spikes to be undetected, which can cause a single unit to be missed. If a threshold is too low, it can allow the detection of noise and non-sortable activities, which would increase the fraction of time spent in the refractory period and, in turn, reduce detection for larger spikes \cite{laboy-juarez_normalized_2019}. Thus, the selected threshold influences the detected events and later spike classification and sorting stages. Therefore, amplitude thresholding detects spikes by their size, which is limited when the spikes vary in amplitude, especially in microneurography, where spikes from C-fibres have a low signal-to-noise ratio \cite{troglio_hybrid_2026}. 
\bigskip

\noindent
An alternative method to amplitude thresholding is template matching, where the spikes are detected by their shape rather than their amplitude. A template of the target waveform correlates with the signal, and detections are the positions of high similarity \cite{laboy-juarez_normalized_2019}. Unlike amplitude thresholding, template matching can detect low-amplitude spikes that some amplitude threshold would miss \cite{laboy-juarez_normalized_2019}, and by normalising the similarity score (normalised cross-correlation), it becomes independent of amplitude, so small or large spikes with the same shape are detected equally. Nonetheless, its performance depends on how distinct the target's spike shape is from other non-targets. 
\bigskip

\noindent
For the spike sorting stage, unsupervised clustering is commonly used, which is where the spikes are grouped by similarity without labels \cite{lewicki_review_1998}. This method is suitable if ground truth labels are not available; however, when the spikes from different units overlap in feature space, clustering does not separate them. The alternative method is used when ground truth labels are available: classification, where the model learns a decision boundary from the labelled target and nontarget spikes \cite{troglio_hybrid_2026}. Before the sorting stage, features have to be extracted to reduce the dimensionality of the waveforms. The most commonly used method for dimensionality reduction is PCA, where the waveforms are projected into the direction of greatest variance \cite{quiroga_unsupervised_2004}. But this can discard discriminative details such as the fine features that can distinguish similar spikes \cite{quiroga_unsupervised_2004}. The alternative is to use the features of the raw waveform directly, similarly to \cite{troglio_hybrid_2026}, who compared feature sets for spike sorting and found that raw waveform features performed better than PCA. However, \cite{troglio_hybrid_2026} compared this for single-channel analysis, so dimensionality reduction could be more useful for multi-channel analysis, where PCA can capture the cross-channel spatial signature. 
\bigskip

\noindent
Separating the target unit by shape alone is challenging as waveforms can be too similar in shape \cite{troglio_hybrid_2026}. Using the marking method and applying a low-frequency electrical stimulus evokes a spike in the target fibre at a consistent latency after each stimulus, and repeated responses align into a distinct vertical in a waterfall plot \cite{schmelz_delayed_1995}. This allows the low-amplitude spikes to be identified by their latency, not just shape \cite{troglio_hybrid_2026}. 
\bigskip

\noindent
The activity-dependent slowing (ADS) is what makes the latency fibre-specific, so when a fibre fires more, its conduction velocity decreases, causing its latency to increase \cite{serra_activitydependent_1999}. Different fibres have different ADS profiles, so their latency tracks diverge, which allows fibres with similar waveforms to be distinguished by how their latency changes over time \cite{troglio_hybrid_2026}. Therefore, for electrically evoked activity, where stimulus timing is available, these knowledge-driven constraints. However, when stimulus timing is not available or is insufficient, for example for spontaneous or chemically evoked activity and for overlapping spikes,  even latency constraints are limited \cite{troglio_hybrid_2026}. 
\bigskip

\noindent
\section{Signal preprocessing and denoising}
\label{sec:denoising}
Denoising these low-signal-to-noise ratio microneurography recordings is an attractive preprocessing step for such analysis. The most common first step for denoising neural recordings is bandpass filtering, which works by isolating specific frequencies that remove the slow components related to the local field potential (LFP), while the upper limit filters out the rapid electrical static that distorts the shape of action potentials \cite{kechris_removing_2021}. However, the standard band-pass filtering introduces sign distortion that alters the waveforms because the filter changes the amplitude and the phase of different frequencies, such as ringing, temporal smearing and baseline shift \cite{de_cheveigne_filters_2019, wiltschko_wavelet_2008}. 
\bigskip

\noindent
To overcome this limitation of bandpass-filtering, wavelet denoising was introduced as it does not force a rigid trade-off between time and frequency resolution, allowing a fine analysis of the signal's frequency components without losing temporal information \cite{baldazzi_systematic_2020}. Yet, because the frequency of neural signals and stimulus artefacts overlap, wavelet struggles to separate noise from signal, because filtering the frequency components of the artefact can lead to the distortion of neural signal \cite{zhang_dual-domain_2026}.  Therefore, researchers are implementing machine learning as a data-driven alternative.  
\bigskip

\noindent
Machine learning can learn the spatial and temporal characteristics of biological signals, and it can capture complex, non-linear relationships between noise and signal \cite{ribeiro_denoising_2023}. Specifically, deep learning approaches such as Autoencoders. Autoencoders learn to reconstruct a clean signal from a noisy input by learning the structure of the data \cite{kechris_removing_2021, bengio_generalized_2013}. These data-driven approaches have shown promising performance in the literature, especially for physiological signal denoising. For example, \cite{lecoq_removing_2021} show improved SNR distribution in 2D neural physiological images using deep encoder-decoder models; \cite{leite_deep_2018} designed a convolutional autoencoder to denoise EEG signals and showed improved results when compared to bandpass approaches; \cite{arsene_design_2020} also show that convolutional networks outperform wavelet methods on denoising ECG. A study by \cite{kechris_removing_2021} used a convolutional autoencoder to denoise extracellular signal to retrieve a signal close to the ground truth, which outperformed wavelet denoising (DWT and SWT). However, the extracellular action potentials were simulated with the MEArea library, and the model was evaluated on simulated data.
\bigskip

\noindent
For spike detection and sorting, the ideal scenario is that these denoisers clean the signal to facilitate a sufficient spike sorting analysis. However, \cite{kay_risk_2022} argues that denoising can introduce systemic bias, where a method can improve the correspondence between the noisy dataset and a ground truth measure, but in turn distort the underlying signal, which may lead to incorrect inference. With spike processing specifically, criteria for noise removal such as reconstruction error can ignore whether the spike morphology is preserved; even if the signal appears cleaner, the spikes may be distorted \cite{baldazzi_systematic_2020}. In addition, autoencoders can show further risk; they can create spurious modes or false attractors in regions that project noise into structured presentation in training \cite{bengio_generalized_2013}. This means for spike detection and sorting, a noise segment that is far from the spike waveforms in training could be reconstructed to be learned as a spike shape or spurious attractor, producing a spike-like waveform that does not originally exist in the raw signal, which could lead to denoiser-induced false detections. So, while denoising can clean the signal, whether it can help downstream detection or sorting is unclear for single-channel. This motivates the evaluation of whether denoising would help spike detection and sorting or would distort the spike and add spurious detections. 
\bigskip

\noindent
\section{Research gap}
\label{sec:research_gap}
The sections above identify three issues. First, Spikes from different fibres have similar waveform shapes to separate for a single-channel, where the spatial information is missed \cite{hill_quality_2011, pedreira_how_2012, troglio_hybrid_2026}. Second, timing based constrains like latency and ADS address this limitation, but do not resolve the waveform shape separation issue for spontaneous and overlapping activity \cite{troglio_hybrid_2026}. Third, denoising is applied in the literature for low-SNR signals, but its benefit for detection and classification is unstated.
\bigskip

\noindent
A further methodological gap is worth flagging: detection and classification evaluations are often susceptible to information leakage, in which the same data is used both to select parameters and to evaluate performance. This double-dipping can inflate apparent performance and lead a classifier to produce optimistic results \cite{kriegeskorte_circular_2009}.
\bigskip

\noindent
This project addresses these gaps by evaluating whether denoising improves signal channel spike detection and classification. A leakage-safe evaluation protocol was applied, in which all tuning is performed on a development set, and performance is scored once on the held-out evaluation set. To assess whether the findings are consistent, the evaluation is performed across multiple channels independently. 
\bigskip

\noindent
Three questions this analysis tries to answer: What is the quality of denoising in terms of noise reduction and spike preservation?; does denoising improve spike detection?; and does denoising improve classification of the target spike against other non-target spikes?. 
% -----------------------------------------------------------------------------
\chapter{Execution}
\label{chap:execution}
This chapter describes the data and the methods used to detect, classify and denoise the target fibre's spikes in a single channel. It first introduces the recording and how the target unit was identified, then estimates the noise floor. It then describes the detection pipeline in three stages: building and matching a spike template, and constraining detections by their timing after the stimulus. It then describes the two classifiers used on the detection survivors: a one-class and two-class support vector machine. Finally describes three denoising methods: Wavelet, Dense Autoencoder and Convolutional Autoencoder to evaluate each method and to assess whether denoising improves detection or classification, and the metrics used throughout. 

\section{The data: rat saphenous nerve microneurography}
This project uses data that consist of time series recordings from the saphenous nerve of a rat, sampled at 30 kHz over approximately 2.3 hours. The nerve carries many afferent fibres, so the electrode records the action potentials of many units at once, superimposed on background noise. Recordings were made using two multi-electrode sensor units placed at different positions along the nerve, each containing several channels corresponding to slightly different electrode locations.
\hspace{0.5cm}

\noindent
The eight channels carry manually labelled ground truth spikes (CH3, CH8, CH9, CH10, CH14, CH19, CH25, CH30). During the recording, an electrical stimulus (Pain stimulus) was delivered repeatedly, giving 2,310 stimulation pulses (TTL) at a normal rate of around 0.25 Hz. The target fibre responds to each stimulus by firing at an approximately fixed latency. The ground truth spike times were marked manually by a human annotator (Figure \ref{fig:raw signal}), who only
labelled recognisable nerve responses; therefore, the annotations are conservative, and not every response is labelled. Each spike is marked at the point where the signal crosses zero on its descending slope. 
\begin{figure}[htbp]
    \centering
    \includegraphics[width=\textwidth]{figures/raw signal with target spikes in blue and stimuli dashed red.png}
    \caption{Segment representing the raw signal of channel 19 with the manually annotated spikes in blue and the stimulus times in red. Each target spike follows a stimulus at a latency of around 100 ms.}
    \label{fig:raw signal}
\end{figure}
\subsection{Target unit and channels}
The lab identified a single target unit as a strong unit with large identifiable spikes on multiple channels. This project analyses five of these channels: 9, 14, 19, 25, and 30, with a strong signal, and each is analysed independently. 
The spike latencies of the target unit were manually annotated on each channel; consequently, they are treated as reference annotations rather than complete ground truth, as they were manually identified and may contain errors. 
Only a 15-minute window was used for analyses, beginning at the 900th stimulus pulse. 
This window length was chosen because the target unit's latency remained stable; a longer window showed activity-dependent slowing that spreads the latency and prevents a fixed latency window from capturing the spikes. The signals were preprocessed by the recording laboratory, so it has been band-pass filtered 300-6000 Hz, common average referenced, and stimulus blanked. 
\section{Noise estimation and preprocessing}
As mentioned, no further filtering was applied in this project, because the recording given was already bandpass-filtered, with stimulation artefact blanked (set to zero). 
The target units are small, only about two to three times the amplitude of the background noise; consequently, they do not stand out clearly and can not be isolated by amplitude threshold alone. Therefore, removing the noise, detecting and classifying the target unit's spikes are the central challenges this chapter tries to address. The noise level is estimated using the median absolute deviation (MAD) \cite{quiroga_unsupervised_2004, park_deep_2019}:
\begin{equation}
    \sigma_{\text{noise}} = \frac{\operatorname{median}(|x|)}{0.6745}
    \label{MAD noise estimate}
\end{equation}

where $x$ is the signal. The median is used instead of standard deviation because the spikes are infrequent, so the median does not inflate the estimate the way standard deviation would; dividing by 0.6745 rescales the value to be comparable to a standard deviation under an assumption of Gaussian noise \cite{drebitz_optimizing_2019}. MAD is an estimate, so it assumes the signal is mostly noise and approximately Gaussian \cite{coupe_robust_2010}. But because it is applied identically to every signal, it is a valid basis for comparison.
\section{Leakage safe evaluation protocol}
For unbiased performance, a chronological split was used on the recording by dividing the recording into a development set on the first 70\% [3602.7-4232.7 s] and a held-out evaluation on the last 30\% of the recording [4232.7-4502.7 s]. All hyperparameter tuning was conducted on the development set, and the evaluation set was scored once to prevent data leakage and to avoid overfitting \cite{rashid_spike_2025}. All components that could carry information across the split were fitted on the development set only; the spike template was built from target spikes in the development annotations split only; the autoencoders were trained on synthetic data derived from the development set. Detection parameters like shape similarity and latency window were selected by maximising F1 on the development set. The classifiers and the two-class classifier were fitted on the development survivors only with fixed hyperparameters and scored once on the evaluation survivors. The confidence intervals were computed by stratified bootstrap resampling of the evaluation candidates.  

\section{Detection}
\subsection{Template matching and spike detection}
The spike template was built by averaging the annotated unit's spikes on the development set, to prevent leakage in later stages. For each marked spike, a 90-sample window (45 samples on either side of the spike trough) is then extracted from the signal and then averaged to give the unit's mean waveform. 
Before averaging, each spike is aligned on a common reference point within a short search region around the marked time, so that the same feature falls at the same position in every window. Without this, the waveform would sit at a slightly different offset, and the average would be blurred. Trough alignment is the standard approach in the literature, though it is sensitive to noise, as the exact location of the trough can shift slightly between spikes \cite{lee_non-linear_2021}. 
The template is slid across the signal to locate spikes that match the unit's waveforms, then a similarity score is computed at every position. Following the normalised correlation approach to spike detection of \cite{laboy-juarez_normalized_2019, hwang_spike_2014}. The score used here is the normalised cross-correlation coefficient between the template and the signal window at each position \cite{article}. Because the mean is subtracted from both the template and the signal window, this is equivalent to the Pearson correlation coefficient: 
\begin{equation}
    r_i = \frac{\sum_{j=1}^{n} (t_j - \bar{t})(x_{i,j} - \bar{x}_i)}
    {\sqrt{\sum_{j=1}^{n} (t_j-\bar{t})^2}\;\sqrt{\sum_{j=1}^{n} (x_{i,j} - \bar{x}_i)^2}}
\end{equation}
Where $t_j$ are the template samples, $x_{i,j}$ are the signal samples in the window at position $i$, $\bar{t}$ and $\bar{x}$ are their mean, and $n$ is the template length. The score $r_i$ lies in [-1,1], and 1 means a perfect match in shape; because the score is normalised, it measures waveform shape independently of amplitude, so a small spike and a large spike with the same shape would score equally.
\bigskip

\noindent
Latency constraints were applied on the detection search space following \cite{troglio_hybrid_2026}. The target fibre fires at around the same latency after each stimulus, which is around 100 ms, so a detection's post-stimulus latency indicates whether it belongs to the target unit or not. For each spike, the latency was computed as the interval from the most recent preceding stimulus; the median latency across the development spikes was around 100 ms. So, a detection is only retained if its post-stimulus latency falls within a window centred on this median. 
Following the leakage safe protocol, both the shape similarity and the latency window width were selected only on the development set by sweeping candidate values and choosing the combination that maximises the development $F_1$ score. These then were frozen and applies to the held out evaluation set. 

\section{Denoising methods}
Three denoisers are used to remove the noise from the recording: wavelet, dense and convolutional autoencoder. 
\subsection{Wavelet}
Wavelet denoising is a signal processing method that works by transforming the signal into wavelet coefficients, removing or shrinking the small coefficients that represent noise, and then reconstructing the signal from the remaining coefficients \cite{wei_wavelet_2005, donoho_minimax_1998}. For this project, SWT was used over DWT for the reasons mentioned in (Chapter \ref{chap:background}). The signal was decomposed using Daubechies 4 to decomposition level 4 to give approximation and detail at each level. Then the noise level was estimated at each level using the same MAD estimator on the detail coefficients at each level. A universal threshold was applied using the VisuShrink formula \cite{donoho_ideal_1994, donoho_-noising_1995}:
\begin{equation}
    \lambda = \sigma \sqrt{2 \ln N}
    \label{VisuShrink}
\end{equation}
Where $\sigma$ is the noise level, estimated using MAD, and $N$ is the number of coefficients at that level. The detail coefficents were shrunk using a soft threshold:
\begin{equation}
    \eta(x, \lambda) = \text{sign}(x)\, \max(|x| - \lambda,\, 0)
    \label{Soft thresholding}
\end{equation}
This then shrinks the coefficients towards zero instead of removing them, and the approximation coefficients were unchanged. The denoised signal is then reconstructed back from the cleaned coefficients to a time-domain signal by inverse SWT. To allow for smoother reconstruction and avoid sharp discontinuities, the soft threshold was chosen over a hard threshold \cite{donoho_-noising_1995,zhang_improved_2019}. 
\subsection{Dense and Convolutional Autoencoders}
Two autoencoders were used to denoise the signal: dense and convolutional autoencoders. Both were trained on the same development set and applied to the evaluation set in the same way; they only differ in their internal architecture. During training, autoencoders generally require clean signals as a target reference so the model can calculate its reconstructive loss \cite{vincent_extracting_2008, kechris_removing_2021}. With real extracellular recordings, there is no clean reference, so training pairs were generated using synthetic pairs from the real spike template. The continuous signal was divided into 60-sample windows; each window was created as one of two types: a window that contains the template spike from the development set or only a noisy window. In the window with a spike, the template was randomly shifted and scaled so the model learns to see spikes at different positions and amplitudes. For each clean window, Gaussian noise was then added to produce the noisy input, giving pairs of noisy input and clean target. Both autoencoders used the Adam optimiser for training to minimise the mean squared error (MSE) between the reconstruction output and the clean target. 
\bigskip

\noindent
To denoise the recording, the trained models were applied to overlapping windows sliding across the signal. Each of the windows was normalised by the MAD noise level before the model to match the training normalisation, then rescaled back, and the outputs were averaged where windows overlapped. 
The differences between Dense and convolutional are that Dense uses fully connected layers that compress the 60-sample window to 30, then to a 15-unit latent space \cite{li_aecuration_2025}. Whereas convolutional replaces the fully connected layers with 1D convolutional layers, where each layer slides a small filter (kernel size 3) across the window to detect local spike-shaped patterns where they occur in the window \cite{kechris_removing_2021, li_accurate_2020}. This handles spikes at different positions better than dense. 
\bigskip

\noindent
The quality of each denoiser was assessed by measuring the noise level of each denoised signal using MAD (Eq\ref{MAD noise estimate}). However, reduction in noise level alone does not indicate a better signal for spike detection or sorting; focusing only on noise removal ignores whether the spike morphology is preserved \cite{baldazzi_systematic_2020}. However, most studies like \cite{baldazzi_systematic_2020} use RMSE to quantify morphology preservation against a clean reference. Here, the recoding is real and noisy annotated spikes with no clear reference to compare against. So, here peak-to-peak was measured following \cite{adam_comparative_2025, harris_accuracy_2000} as a quantitative waveform feature to assess the spike amplitude after denoising. Peak-to-peak evaluates the denoisers by quantifying the absolute difference between the maximum and minimum values of a waveform before and after filtering \cite{harris_accuracy_2000}. For each of the target spikes, the peak-to-peak amplitude was computed from the aligned waveforms. Amplitude retention was then expressed as the mean peak-to-peak amplitude of each denoised signal relative to the raw signal. 

\section{Classification}
Two support vector machine (SVM) classifiers with radial basis function (RBF) kernels were applied to the development set through aligned waveforms, following the raw waveforms approach of \cite{troglio_hybrid_2026}, which found that raw waveforms performed better than PCA-derived features. One-class and two-class SVMs were used for this analysis. The difference between the two SVMs: one-class was trained on only the target survivors of the development set, and two-class was trained on both target and non-target survivors of the development set. Both SVMs were trained with fixed hyperparameters (nu = 0.1, C = 1.0). Then, each fitted classifier with its feature scaler was applied once on the held-out evaluation set survivors. This ensures that no candidate is scored by a model trained on it, which avoids data leakage. 
\subsection{One-class SVM}
One-class SVM is a novelty detection machine learning method that learns a decision function to classify raw data as similar to or different to the training set \cite{yang_efficient_2021}. The one-class model was trained on the target survivors in the development set waveforms and then evaluated once on the held-out evaluation set target and non-target survivors. Because the model was trained on the development set and scored only on the evaluation set, no waveform appears in both training and testing. So, each evaluation candidate received a single prediction from a fixed model. An RBF was used ($\gamma$ = scale), with nu ($\nu$) fixed at 0.1 a priori to avoid tuning the hyperparameter on the evaluation data. This prevents biased estimates of how well the model would generalise, because nu tests the expected fraction of training points allowed outside the learned boundary and the fraction of support vectors in the interval [0, 1]; and gamma controls the influence of each training example \cite{moctezuma_multi-objective_2020}. 
\bigskip

\noindent
Studies such as \cite{troglio_hybrid_2026} use the $F_1$ score as a primary measure of the one-class model. However, $F_1$ relies on binary predictions that are evaluated at a fixed decision boundary, so they are sensitive to the boundary tuning hyperparameter. The nu hyperparameter in one-class controls the upper bound on the fraction of outliers and the lower bound on support vectors, influencing how tightly the boundary shrinks or how widely it expands around the data \cite{scholkopf_estimating_2001}. According to \cite{tax_support_2004}, altering this boundary dictates the trade-off between false positives and false negatives, so if nu does not exactly match the true contamination rate, the target precision-recall balance would shift and cause the discrete $F_1$ score to fluctuate. 
\bigskip

\noindent
For this analysis, these metrics were used, but not as the primary measurement for SVM performance; instead, ROC-AUC is used with a bootstrap confidence interval. This is because ROC-AUC evaluates the continuous decision output across all thresholds, so adjusting nu does shift the decision boundary, but it preserves the candidate ranking order, making ROC-AUC invariant to nu tuning \cite{tax_support_2004, abdelhamid_balancing_2024}. 
Because one-class RBF classifiers are sensitive to feature scaling \cite{juszczak_feature_2002, feragen_adaptive_2015}, the one-class SVM features were not standardised. Standardising these waveforms divides low-variance tail points by near-zero standard deviation. This amplifies background noise, corrupts the waveform's morphology and inflates Euclidean distance in the RBF kernel, which causes the boundary to reject almost all candidates, including target ones. 
\subsection{Two-class SVM}
Unlike one-class SVM, two-class SVM is a supervised machine learning model that works by separating data points into two distinct groups to find a boundary between them \cite{krebs_applying_2024}. So, here it's trained on both target and non-target survivors' waveforms of the development set to learn the boundary the separate these two classes. Compared to the non-target survivors, the target survivors are a minority, so class imbalance was handled by class weighting rather than resampling. For the two classes, the features were standardised using a StandardScaler fit on the development set and then applied to the evaluation set, in contrast to one-class SVM. 
\bigskip

\noindent
Two-class SVM was standardised following the recommendation of \cite{carugo_users_2010}. Standardisation here is safe because the two-class training set is heterogeneous, containing both target and non-target spikes and diverse background noise waveforms, so the feature standard deviation remains non-zero across all points, eliminating the risk of near-zero division. This in turn allows for balanced feature scales to be standardised and improves the distance between the decision boundary and support vectors without amplifying baseline noise or distorting kernel distance metrics \cite{cortes_support-vector_1995}. 
\bigskip

\noindent
As with one-class, the C parameter was fixed at 1.0 a priori to avoid tuning on the evaluation set. The classifier was trained on the development survivors and applied once to evaluation survivors, so that each evaluation candidate is predicted once from one fixed model and the performance is then measured by ROC-AUC, precision, recall, and F1. 

\section{Evaluation metrics}
 To evaluate detection and classification performance on raw and denoised signals, metrics derived from the confusion matrices were used, along with ROC-AUC with bootstrap confidence intervals for classification. 
 \subsection{Confusion matrix}
 The confusion matrices count true positives (TP), false positives (FP), false negatives (FN), and true negatives (TN) \cite{singh_diagnosing_2021}. From these, precision [TP/(TP+FP)], recall [TP(TP+FN)] and their harmonic mean (the $F_1$ score) were computed. The false-positive-to-true-positive ratio (FP/TP) was also computed for detection following \cite{troglio_hybrid_2026} as an indicator of the false-positive burden. 
 \subsection{ROC-AUC}
 To evaluate the classifiers' performance, the area under the ROC curve (AUC) was used as an indicator of the classifiers' discriminatory power and the model's ability to separate classes across all decision thresholds \cite{hanley_meaning_1982}. ROC-AUC has a fixed chance level of 0.5, so an AUC of 0.5 means no discrimination between classes (equivalent to random guessing), and a higher AUC means the stronger the discriminatory power \cite{hosmer_applied_2013}. As mentioned above, because ROC-AUC is invariant to the class imbalance between the minority target and majority non-target, it was chosen as the primary metric for the classifiers. It is fixed chance reference makes it straightforward to assess whether a classifier performs above chance. 
 \subsection{Bootstrap confidence intervals}
 Each ROC-AUC is a point estimate from one candidate set. The uncertainty is quantified with 95\% confidence intervals computed by bootstrap resampling \cite{efron_resampling_2001}, as used by \cite{troglio_hybrid_2026} for spike sorting performance. The difference is that \cite{troglio_hybrid_2026} quantified performance uncertainty around $F_1$ score, while here around ROC-AUC. The evaluation candidates were resamples with replacement 2000 times, and the ROC-AUC was computed for each resample. Then the 95\% intervals were taken as the 2.5th and 97.5th percentile of the resample values. Any resample that only contained one class was skipped. To interpret the results, a narrow interval indicates AUC is likely accurate, while a wide interval indicates the AUC values are less reliable \cite{corbacioglu_receiver_2023}. 
% -----------------------------------------------------------------------------
\chapter{Critical Evaluation}
\label{chap:evaluation}
\section{Data and channels}
The evaluation was performed on the five channels where the target unit has a strong signal (channels 9, 14, 19, 25, and 30). Each channel is analysed independently as a single-channel recording over 15 minutes window for consistency analysis. Table (\ref{tab:channel data overview}) shows how many target annotations were in each channel and how many were in the 15-minute window analysed; each channel was split chronologically into a development set used for all tuning and a held-out set that was scored once. The noise floor MAD is consistent between channels. 
\begin{table}[htbp]
    \centering
    \begin{tabular}{lccccc}
    Channel & Total annotations & Window annotations & Developemnt & Evaluation & MAD \\
    \toprule
    9 & 1628 & 170 & 108 & 62 & 4.80 \\
    14 & 1639 & 173 & 113 & 60 & 4.88 \\
    19 & 1623 & 177 & 116 & 61 & 4.96 \\
    25 & 1777 & 186 & 120 & 66 & 4.63 \\
    30 & 1414 & 174 & 133 & 61 & 4.82 \\
    \bottomrule
    \end{tabular}
    \caption{Overview of the five channels analysed for this project. The target unit appears with a strong signal on all channels. Each channel is analysed independently over a 15-minute window with a chronological split of 70\% development and 30\% evaluation. MAD is the raw noise level estimate for each channel.}
    \label{tab:channel data overview}
\end{table}
\section{Denoising quality}
The three denoisers were assessed across the five channels using two measures: MAD to assess each denoiser's noise reduction relative to the raw and the preservation of spike amplitude by using peak-to-peak amplitude relative to the raw signal. These measures were used instead of reconstructive error because no clean signal exists for these real recordings to compare against \cite{baldazzi_systematic_2020, adam_comparative_2025}
All denoisers reduced the noise level, as seen in (Table \ref{tab:denoising quality}). Dense reduced the most noise at around 0.50 MAD, and the wavelet denoiser removed the least noise at around 2.30 MAD in comparison to the raw measurement of around 4.80 MAD. Despite the Dense autoencoder achieving the lowest MAD, it only retained 17-20\% of the raw spike amplitude across channels, making it the most destructive. The convolutional autoencoder reduced the noise and preserved the spike amplitude, achieving 55-57\% (Table \ref{tab:denoising quality}; Figure \ref{fig:denosiing quality ch19}). This pattern was consistent across the five channels. 

\begin{table}[htbp]
    \centering
    \begin{tabular}{llccc}
    \toprule
    Channel & Signal & MAD & Peak-to-peak & Retained \\
    \midrule
    \multirow{4}{*}{9} & Raw & 4.80 & 22.99 & 100\% \\
    & Wavelet & 2.23 & 7.36 & 31\% \\
    & Dense & 0.48 & 4.47 & 19\% \\
    & Conv & 1.08 & 13.13 & 55\% \\
    \midrule
\multirow{4}{*}{14} & Raw & 4.89 & 23.82 & 100\% \\
    & Wavelet & 2.23 & 7.36 & 31\% \\
    & Dense & 0.48 & 4.47 & 19\% \\
    & Conv & 1.08 & 13.13 & 55\% \\
\midrule
\multirow{4}{*}{19} & Raw & 4.96 & 24.08 & 100\% \\
    & Wavelet & 2.25 & 7.62 & 32\% \\
    & Dense & 0.58 & 4.81 & 20\% \\
    & Conv & 1.25 & 13.67 & 57\% \\
\midrule
\multirow{4}{*}{25} & Raw & 4.63 & 21.75 & 100\% \\
    & Wavelet & 2.31 & 7.77 & 36\% \\
    & Dense & 0.50 & 3.91 & 18\% \\
    & Conv & 1.05 & 11.56 & 53\% \\
\midrule
\multirow{4}{*}{30} & Raw & 4.82 & 22.75 & 100\% \\
    & Wavelet & 2.25 & 7.22 & 32\% \\
    & Dense & 0.48 & 3.92 & 17\% \\
    & Conv & 1.23 & 12.68 & 56\% \\
    \bottomrule
    \end{tabular}
    \caption{Denoising quality across the five channels. Each denoiser reduces the noise floor (MAD) and preserves the spike amplitude peak-to-peak relative to raw. The dense autoencoder shows the lowest MAD and retains the least amplitude across all channels, while the convolutional autoencoder reduces the noise floor but preserves the most spike amplitude across all channels.}
    \label{tab:denoising quality}
\end{table}
\begin{figure}[H]
    \centering
    \includegraphics[width=1.1\textwidth]{figures/denoising_quality_leakage_safe.png}
    \caption{Effect of each denoiser on the signal on channel 19. Left is an overview of a segment around a target spike. Right is a detail view of the spike.}
    \label{fig:denosiing quality ch19}
\end{figure}

This suggests that MAD alone is not enough to determine denoising quality, because a distorted spike carries less information for detection and classification. So, only focusing on noise removal ignores whether the spike morphology is preserved \cite{baldazzi_systematic_2020}. However, \cite{baldazzi_systematic_2020} used RMSE to quantify morphology preservation against a clean reference. Here, the recording is real and noisy with annotated ground truth, so peak-to-peak was measured following \cite{adam_comparative_2025}. The results for denoising quality are consistent across the five channels; the raw and convolutional spikes retain their depth, while the wavelet and dense did not. 

\section{Template matching detection}
This section compares detection before and after denoising, comparing the three denoisers across the five channels. 
\subsection{Baseline detection}
 Detection was applied to the raw signal on each channel and scored against the reference annotations in the held-out evaluation period (Table \ref{tab:Baseline detection}; Figure \ref{fig:detection baseline}). The false positive to true positives ratio (FP/TP) across all channels exceeded 1, ranging from 2.75 to 4.73, and recall was low, around 0.18 to 0.27. Detection on the raw signal of every channel showed more false positives than true positives. This indicates that these false positives resemble the target spikes in both waveforms and latency window, as they passed both the shape similarity threshold and the latency window. These results are consistent with the literature. According to \cite{troglio_hybrid_2026, hill_quality_2011} in single-electrode microneurography analysis, waveform shapes are highly similar across fibres, and it is difficult to separate overlapping units on a single contact. The FP/TP ratio here for the raw signal is around 2.75 to 4.73, which is lower than those reported by \cite{troglio_hybrid_2026} for multi-fibre recordings; the FP/TP is around 2 to 20. However, these studies did not denoise the signal; the next subsection evaluates whether denoising the signal would make spike detection by shape and latency any better. 

 \begin{table}[htbp]
     \centering
     \begin{tabular}{lcccccccc}
     \toprule
    Channel & Survivors & TP & FP  & FN & Precision & Recall & F1 & FP/TP \\
    \midrule
    9  & 61 & 12 & 49 & 50 & 0.197 & 0.194 & 0.195 & 4.08 \\
    14 & 60 & 16 & 44 & 44 & 0.267 & 0.267 & 0.267 & 2.75 \\
    19 & 63 & 11 & 52 & 50 & 0.175 & 0.180 & 0.177 & 4.73 \\
    25 & 58 & 14 & 44 & 52 & 0.241 & 0.212 & 0.226 & 3.14 \\
    30 & 62 & 13 & 49 & 48 & 0.210 & 0.213 & 0.211 & 3.77 \\
    \bottomrule 
     \end{tabular}
     \caption{Per-channel baseline detection on raw signal on held-out evaluation. The false positives to true positive ratio (FP/TP exceeds 1 for every channel, with channel 19 the highest and channel 14 the lowest.}
     \label{tab:Baseline detection}
 \end{table}
\begin{figure}[H]
     \centering
     \includegraphics[width=\textwidth]{figures/Detection performamce befoer denoiisng (1).png}
     \caption{Baseline detection performance per channel on the raw signal. Spike detection performance using similarity thresholding and a latency window: TP, FP, and FN are visualised on a log scale for the raw signal across the five channels on the held-out evaluation period. All channels have higher false positives compared to the true positives, with Ch19 exhibiting the most false positives and the least true positives.}
     \label{fig:detection baseline}
 \end{figure}
\subsection{Detection after denoising}
Detection was then applied to each denoised signal from wavelet, dense and convolutional denoisers and compared to the raw baseline (Table \ref{tab:detection denoised}; Figure \ref{fig:detection after denoising}. 
The effect of denoising on the FP/TP is inconsistent across channels, and no denoiser improved the FP/TP; the ratio remains above 1 for all denoisers and all channels. 
Wavelet improved the ratio slightly on three channels (ch9, ch14, and ch19), and worsened the ratio on two channels (ch25 and ch30); dense also improved three (ch9, ch25, and ch30) and worsened the ratio on two channels (ch14, 19); convlutional improved 2 (ch9, ch19) and showed worse ratio on (ch14, ch25, ch30). Only channel 9: the three denoisers reduced FP/TP, although even there the ratio remained above 2.7.

\begin{table}[htbp]
    \centering
    \begin{tabular}{llccccccc}
    \toprule
    Channel & Signal & Survivors & TP & FP & FN & Recall & F1 & FP/TP \\
    \midrule
    \multirow{4}{*}{9}  
    & Raw & 61 & 12 & 49 & 50 & 0.194 & 0.195 & 4.08 \\
    & Wavelet & 59 & 16 & 43 & 46 & 0.258 & 0.264 & 2.69 \\
    & Dense & 60 & 15 & 45 & 47 & 0.242 & 0.246 & 3.00 \\
    & Conv & 65 & 16 & 49 & 46 & 0.258 & 0.252 & 3.06 \\
    \midrule
    \multirow{4}{*}{14} 
    & Raw & 60 & 16 & 44 & 44 & 0.267 & 0.267 & 2.75 \\
    & Wavelet & 59 & 16 & 43 & 44 & 0.267 & 0.269 & 2.69 \\
    & Dense & 55 & 10 & 45 & 50 & 0.167 & 0.174 & 4.50 \\
    & Conv & 65 & 15 & 50 & 45 & 0.250 & 0.240 & 3.33 \\
    \midrule
    \multirow{4}{*}{19} 
    & Raw & 63 & 11 & 52 & 50 & 0.180 & 0.177 & 4.73 \\
    & Wavelet & 58 & 14 & 44 & 47 & 0.230 & 0.235 & 3.14 \\
    & Dense & 64 & 8 & 56 & 53 & 0.131 & 0.128 & 7.00 \\
    & Conv & 63 & 16 & 47 & 45 & 0.262 & 0.258 & 2.94 \\
    \midrule
    \multirow{4}{*}{25} 
    & Raw & 58 & 14 & 44 & 52 & 0.212 & 0.226 & 3.14 \\
    & Wavelet & 52 & 12 & 40 & 54 & 0.182 & 0.203 & 3.33 \\
    & Dense & 61 & 21 & 40 & 45 & 0.318 & 0.331 & 1.90 \\
    & Conv & 61 & 12 & 49 & 54 & 0.182 & 0.189 & 4.08 \\
    \midrule
    \multirow{4}{*}{30} 
    & Raw & 62 & 13 & 49 & 48 & 0.213 & 0.211 & 3.77 \\
    & Wavelet & 51 & 7 & 44 & 54 & 0.115 & 0.125 & 6.29 \\
    & Dense & 65 & 14 & 51 & 47 & 0.230 & 0.222 & 3.64 \\
    & Conv & 62 & 11 & 51 & 50 & 0.180 & 0.179 & 4.64 \\
    \bottomrule
    \end{tabular}
    \caption{Per channel and signal detection after denoising on the held-out evaluation. Denoising did not reduce the FP/TP ratio; all are exceeding 1 on every channel, and it shows an inconsistent effect across channels.}
    \label{tab:detection denoised}
\end{table}
\begin{figure}[htbp]
     \centering
     \includegraphics[width=\textwidth]{figures/Detection performamce after denoisng.png}
     \caption{Detection performance per channel before and after denoising the signal. The TP, FP, and FN are visualised on a log scale for the denoised signal across the five channels on the held-out evaluation period. Denoising does not consistently reduce the false positive count.}
     \label{fig:detection after denoising}
 \end{figure}

When comparing spike template matching detection before denoising to after denoising, none of the three denoisers reduced the FP/TP ratio to below 1 on any channel; even the largest improvement, which was dense on CH25, only improved it to 1.90 compared to the raw ratio of 3.14. The scattered pattern of improvement and worsening is quite balanced across the fifteen denoiser channel combinations. 
\bigskip

\noindent
Several reasons and mechanisms reported in the literature for denoising increasing the number of false positives with some denoisers and some channels, and these fall into two categories. These false positives could be genuine biological spikes that were buried by the noise and were detected by the template matching after denoising, or artefacts introduced by the denoisers themselves.
According to \cite{bengio_generalized_2013}, denoising autoencoders can learn to pull corrupted points back onto the manifold via a learned vector field, so when the model receives unstructured noise, it can reshape that noise into the nearest pattern, which, in this case, is spikes. This could be the reason for the higher number of false positives compared to raw. Furthermore, denoising can diminish the spikes' fine features that help distinguish them from neighbouring units \cite{baldazzi_systematic_2020}, which in turn causes non-target spikes to resemble the target template. 
\bigskip

\noindent
If the additional detections (false positives) by the denoisers are actual spikes, they should look like the spike's waveforms in the raw signal; and if they are artefacts, they should be absent or near flat in the raw signal since they were created. This check was only done on ch19 for the dense autoencoder, where the FP/TP ratio increased from 4.73 to 7.00 with 11 additional detections. These additional detections have a mean trough of -11.77 in the raw signal, similar to the target spikes, indicating that these extra false positives might be spike-like events rather than artefacts produced by the denoiser. However, this can not rule out any other possibilities, as only the mean trough spikes were checked and only on one channel with one denoiser. Nonetheless, in a single channel, spikes do overlap, have similar waveform shapes and latency, and without the spatial information that of multi-channel analysis, it is difficult to distinguish target spikes from non-target \cite{troglio_hybrid_2026}.
\section{Classification}
\subsection{Baseline classification}
The one-class and two-class SVMs were applied to the survivors from detection on the raw signal, trained on the development set and scored on the held-out evaluation set once.  Both (Table \ref{tab:classification baseline}; Figure \ref{fig:classification_baseline}) show the ROC-AUC with 95\% confidence interval included 0.5. For one class svm the ROC-AUC for all channels ranged from 0.455 to 0.602, while the two-class SVM ranged from 0.399 to 0.606, and all confidence intervals included 0.5. The confusion matrix counts (Table \ref{tab:classification baseline confusion metrics}) show that for every channel, the false positives are more than true positives and the precision is between 0.15 and 0.30. 
\begin{table}[H]
    \centering
    \begin{tabular}{lcc}
    \toprule
    Channel &  One-class ROC-AUC [95\% CI] & Two-class ROC-AUC [95\% CI] \\
    \midrule
    9 & 0.602 [0.446, 0.755] & 0.481 [0.299, 0.658] \\
    14 & 0.477 [0.303, 0.651] & 0.506 [0.347, 0.666] \\
    19 & 0.455 [0.264, 0.652] & 0.502 [0.283, 0.719] \\
    25 & 0.539 [0.370, 0.711] & 0.399 [0.237, 0.571] \\
    30 & 0.523 [0.356, 0.684] & 0.606 [0.411, 0.780] \\
    \bottomrule
    \end{tabular}
    \caption{Per-channel baseline classification performance on the raw signal. All ROC-AUC confidence intervals are wide, but all include 0.5, indicating that the model's performance on the raw signal is indistinguishable from chance.}
    \label{tab:classification baseline}
\end{table}
\begin{table}[H]
    \centering
    \begin{tabular}{llccccccc}
    Classifier & Channel & TP & FP & FN & TN & Precision & Recall & F1 \\
    \midrule
    \multirow{5}{*}{One-class}
    & 9  & 6  & 19 & 6 & 30 & 0.240 & 0.500 & 0.324 \\
    & 14 & 10 & 28 & 6 & 16 & 0.263 & 0.625 & 0.370 \\
    & 19 & 6  & 30 & 5 & 22 & 0.167 & 0.545 & 0.255 \\
    & 25 & 8  & 19 & 6 & 25 & 0.296 & 0.571 & 0.390 \\
     & 30 & 8  & 28 & 5 & 21 & 0.222 & 0.615 & 0.327 \\
    \midrule
    \multirow{5}{*}{Two-class}
     & 9  & 7  & 28 & 5 & 21 & 0.200 & 0.583 & 0.298 \\
     & 14 & 12 & 33 & 4 & 11 & 0.267 & 0.750 & 0.393 \\
     & 19 & 6  & 33 & 5 & 19 & 0.154 & 0.545 & 0.240 \\
     & 25 & 7  & 26 & 7 & 18 & 0.212 & 0.500 & 0.298 \\
     & 30 & 8  & 22 & 5 & 27 & 0.267 & 0.615 & 0.372 \\
    \bottomrule
    \end{tabular}
    \caption{Per-channel confusion matrix counts of baseline classification. The false positives exceed the true positives on every channel.}
    \label{tab:classification baseline confusion metrics}
\end{table}
\begin{figure}[H]
    \centering
    \includegraphics[width=\textwidth]{figures/classification_baseline.png}
    \caption{Dot plot showing per-channel ROC-AUC scores for baseline classification on raw signal. The ROC-AUC is shown with 95\% bootstrap confidence intervals for one class on the left side and two classes on the right side. Every interval includes the chance level of 0.5 where the red dashed line.}
    \label{fig:classification_baseline}
\end{figure}
\subsection{Classification after denoising}
On denoised signals, no consistent improvement in classification performance was observed across channels (Table \ref{tab:classificatio-roc denoised}; Figure \ref{fig:classification roc after denoising}). The biggest improvement observed with one-class SVM after denoising the signal was for ch30, where ROC-AUC increased from 0.523 (95\% CI of 0.356-0.684) to 0.675 (95\% CI of 0.409-0.896) following wavelet denoising and 0.649 (95\% CI of 0.483-0.791) following conv denoising. An improvement was also observed on ch14 with dense denoising, where ROC-AUC increased from 0.477 (95\% CI of 0.303-0.651) to 0.589(95\% CI of 0.411-0.764). In contrast, the ch19 denoising signal showed decreased performance, with ROC-AUC decreasing from 0.455 (95\% CI of 0.264-0.652) to 0.370(95\% CI of 0.215-0.522) with the conv denoiser. on the other hand, wavelet reduced two-class performance for ch30 from ROC-AUC 0.606(95\% CI of 0.411-0.780) to 0.344(95\% CI of 0.133-0.578). The full confusion matrix after denoising is reported in the Appendix (Table \ref{tab:confusion matrix classification after denoising}). 
\begin{table}[H]
    \centering
    \begin{tabular}{llcc}
    \toprule
    Channel & Denoiser & One-class ROC-AUC [95\% CI] & Two-class ROC-AUC [95\% CI] \\
    \midrule
    \multirow{3}{*}{9}  
    & Wavelet & 0.541 [0.366, 0.711] & 0.509 [0.331, 0.676] \\
    & Dense & 0.536 [0.370, 0.702] & 0.527 [0.351, 0.704] \\
    & Conv & 0.537 [0.363, 0.703] & 0.509 [0.356, 0.672] \\
    \midrule
    \multirow{3}{*}{14} 
    & Wavelet & 0.461 [0.305, 0.625] & 0.491 [0.308, 0.673] \\
    & Dense & 0.589 [0.411, 0.764] & 0.613 [0.387, 0.822] \\
    & Conv & 0.517 [0.367, 0.680] & 0.585 [0.419, 0.740] \\
    \midrule
    \multirow{3}{*}{19} 
    & Wavelet & 0.435 [0.258, 0.612] & 0.484 [0.292, 0.680] \\
    & Dense & 0.438 [0.239, 0.623] & 0.531 [0.310, 0.757] \\
    & Conv & 0.370 [0.215, 0.522] & 0.387 [0.243, 0.542] \\
    \midrule
    \multirow{3}{*}{25} 
    & Wavelet & 0.552 [0.371, 0.727] & 0.348 [0.192, 0.521] \\
    & Dense & 0.576 [0.421, 0.720] & 0.380 [0.240, 0.523] \\
    & Conv & 0.580 [0.427, 0.735] & 0.587 [0.422, 0.747] \\
    \midrule
    \multirow{3}{*}{30} 
    & Wavelet & 0.675 [0.409, 0.896] & 0.344 [0.133, 0.578] \\
    & Dense & 0.598 [0.420, 0.758] & 0.410 [0.267, 0.564] \\
    & Conv & 0.649 [0.483, 0.791] & 0.613 [0.415, 0.795] \\
    \bottomrule
    \end{tabular}
    \caption {Per-channel classification performance on denoised signal. The ROC-AUC confidence interval is wide, but all include 0.5; denoised signal does not improve classification to above chance.}
    \label{tab:classificatio-roc denoised}
\end{table}
\begin{figure}[H]
    \centering
    \includegraphics[width=\textwidth]{figures/classification after denoising.png}
    \caption{Per-channel and denoiser (wavelet, dense, and convolutional autoencoders) ROC-AUC score for classification performance on denoised signal.The ROC-AUC is shown with 95\% bootstrap confidence intervals for one class on the left side and two classes on the right side. Denoising did not raise performance above chance; every interval includes the chance level of 0.5, where the red dashed line}
    \label{fig:classification roc after denoising}
\end{figure}
Across all denoisers and channels, the confidence intervals are wide, spanning from 0.2 to 0.8. These intervals were computed using stratified bootstrap resampling, so this width is not due to shifting class balance across resamples. However, two factors account for this wide interval: first, the small evaluation set, which was around 60 survivors per channel, of which only around 14 to 22 are target, limiting the statistical power and widening the intervals. The second factor is that the target and non-target waveforms overlap, making it difficult for the classifier to separate them even on a denoised signal, because non-target events appear to have a spike shape similar to the target \cite{troglio_hybrid_2026, hill_quality_2011}. The confusion matrix (Table \ref{tab:confusion matrix classification after denoising}) supports this; the classifier accepts non-target survivors as targets, returning precision [0.13 to 0.40]; false positives exceed true positives, which confirms that classification is not discriminating between target spikes and non-target spikes and that this is consistent between channels. 
These results suggest that the raw waveform features and the SVM tested here are unable to separate target units from non-target spikes on a single channel. This is a negative finding, suggesting that waveform shape alone can not resolve this contamination. Denoising the signal does not change this conclusion. Although some point estimates differ slightly, none of these differences seems meaningful. For example, the two-class ROC-AUC for the conv denoiser reached 0.61 on channels 14 and 30; their confidence intervals still include 0.5 (0.42-0.82) and (0.42-0.80); does not indicate above-chance performance. Similarly, the channel-to-channel variation observed is expected because of the small evaluation set, resulting in ROC-AUC estimates that are quite uncertain, so a difference such as 0.44 vs 0.68 could just be sampling variability rather than evidence that some channels provide better class separation. Furthermore, because multiple channels, denoising and classification conditions are being evaluated, some ROC-AUC may rise due to sampling variability and should therefore be interpreted with caution \cite{sullivan_facts_2021}. 
\section{Summary of the results}
The overall detection performance showed variation across conditions (Table \ref{tab:Summary}). The FP/TP ratio for raw was 3.69 $\pm$ 0.70, compared with 3.63 $\pm$ 1.35 for wavelet, 4.01 $\pm$ 1.72 for dense and 3.61 $\pm$ 0.65 for conv. Consequently, denoising does not give a consistent reduction in the false positive to true positive ratio. The lowest ratio is observed with conv, while the highest is with dense; however, the differences are small relative to the variability across channels. 
Because both classifiers were performed on the survivors retained from detection, changes at the detection stage give important context for the classification results. 

Similarly, the overall performance of classification across the five channels showed that denoising does not produce consistent improvement in classification performance (Table \ref{tab:Summary}). Dense denoising for one-class classification shows the highest mean ROC-AUC (0.547 $\pm$ 0.059), compared with 0.518 $\pm$ 0.051 for the raw condition. In contrast, for two-class SVM, the conv denoising showed the highest mean ROC-AUC (0.536 $\pm$ 0.082), compared to 0.499 $\pm$ 0.066 for the raw condition, and Wavelet showed lower mean ROC-AUC (0.435 $\pm$ 0.073). However, the effect of denoising on classification differed between channels and denoisers (Figure \ref{fig:classification heatmap}). For example, one-class ROC-AUC on channel 30 increased from 0.52 in the raw condition to 0.68 with wavelet-denoised signal and 0.65 with conv. Whereas the ROC-AUC on channel 19 decreased to 0.37 with conv denoising. Similarly, for two-class SVM, conv increased ROC-AUC on channel 25 from 0.40 to 0.59, and reduced to 0.35 with wavelet. 
\begin{table}[H]
    \centering
    \begin{tabular}{lccc}
    \toprule
    Condition & One-class ROC-AUC & Two-class ROC-AUC & Detection FP/TP \\
    \midrule
    Baseline (raw) & 0.518 $\pm$ 0.051 & 0.499 $\pm$ 0.066 & 3.69 $\pm$ 0.70 \\
    Wavelet & 0.533 $\pm$ 0.084 & 0.435 $\pm$ 0.073 & 3.63 $\pm$ 1.35 \\
    Dense & 0.547 $\pm$ 0.059 & 0.492 $\pm$ 0.086 & 4.01 $\pm$ 1.72 \\
    Conv & 0.531 $\pm$ 0.092 & 0.536 $\pm$ 0.082 & 3.61 $\pm$ 0.65 \\
    \bottomrule
    \end{tabular}
    \caption{Summary table of mean $\pm$ standard deviation (SD) of one-class and two-class SVM ROC-AUC and detection FP/TP ratio across the five channels for each condition.}
    \label{tab:Summary}
\end{table}
\begin{figure}[htbp]
    \centering
    \includegraphics[width=0.75\textwidth]{figures/classification_heatmap.png}
    \caption{Classification ROC-AUC across five channels and all conditions (one-class and two-class SVM, before and after denoising). Values of 0.5 are chance-level discrimination. The heatmap displays point estimates only; the associated 95\% confidence intervals are reported in (Table \ref{tab:classificatio-roc denoised})}
    \label{fig:classification heatmap}
\end{figure}
Overall, despite multiple higher point estimates at individual channels and denoising combinations, the mean ROC-AUC values remained close to the chance level of 0.5. Therefore, these higher point estimates should be interpreted carefully because each channel was evaluated using a small number of targets, resulting in uncertainty around the AUC estimates. 
 
% -----------------------------------------------------------------------------

\section{Limitations}
Several limitations must be considered when interpreting the results. The fixed latency window used here assumes the latency is constant, but over longer recordings, activity-dependent slowing (ADS) spreads the latency. Here the method is restricted to a 15-minute window where the latency is stable; longer and more active periods would require latency tracking \cite{troglio_hybrid_2026}. This leads to the next limitation: because of this restriction of 15 minutes window, the held- out evaluation set is small ($\approx$ 60 survivors per channel, of which only $\approx$ 14--22 are targets). This limits the statistical power and gives wide confidence intervals that are less informative. The at-chance findings are consistent with chance rather than proven at chance, so a larger evaluation set would tighten the intervals and would give better interpretability. 
\bigskip

\noindent
This study only tests three denoisers (Wavelet SWT, dense autoencoder and convolutional autoencoder). Other denoising methods such as other wavelet families, other architectures, and self-supervised methods may perform differently. So, the conclusion that denoising does not help is limited to these three denoisers and implementations. 
The SVM hyperparameters in this study were fixed a priori, not tuned, so tuning the hyperparameters may change the performance. 
The target spikes are manual annotations, not perfect ground truth, so some detections scored as false positives might include unlabelled target spikes, which would affect the evaluation. 
Because the data used for the analysis is real, denoising quality was assessed by amplitude preservation and noise level reduction relative to raw, so no reconstruction error RMSE against clean signal. Therefore, this is a practical proxy, not the ideal measure. 
\bigskip

\noindent
Only the raw waveform features were used following \cite{troglio_hybrid_2026}, which found that raw feature waveforms performed better than PCA. However, for this data, a different feature representation might change the classification performance. The findings of this analysis are from one recording with a specific signal-to-noise ratio (SNR) and fibre count condition; different recordings with higher SNR and fewer fibres might show different results, as \cite{troglio_hybrid_2026} found variation across recordings. 
\bigskip

\noindent
Finally, this thesis performed binary classification, which only distinguishes the target spikes from non-target rather than full spike sorting. Full spike sorting would cluster all detected spikes into multiple distinct units, assigning each spike to its fibre. The binary classification used here is therefore a simpler task. 
\chapter{Conclusion}
\label{chap:conclusion}
This thesis evaluated whether denoising improves single-channel spike detection and classification in rat saphenous nerve microneurography. A leakage-safe pipeline was developed: all parameters for the spike template, denoisers and detection were fitted on a development set, and performance was scored once on held-out evaluation data. Three denoisers were compared: wavelet, a dense autoencoder and a convolutional autoencoder and analysis was performed for five independent channels on which the target unit appears.  
\bigskip

\noindent
The main findings of this dissertation: first, all denoisers reduced the noise level but attenuated the spikes. The dense autoencoder reduced the noise the lowest, and the level of spike amplitude preservation was different across denoisers, dense being the most destructive, only retaining 18\% of spike amplitude relative to raw. Second, denoisers did not improve spike detection; the false positive to true positive ratio remained above 1 on every channel. Third, both classifiers were at chance on every channel, before and after denoising, indicating that single-channel waveform shape can not be separated from non-target spike waveform shape. These findings were consistent across all channels.
\bigskip

\noindent
The single-channel leakage-safe pipeline developed for this thesis is complete and functional; it detected spikes, applied three denoisers, two classifiers were applied on the survivors, and the performance was evaluated on the held-out evaluation data, and it was applied to all five channels independently. 
The multi-channel analysis that was planned as the next step, which would combine the channels to exploit the spatial signature of each spike across electrodes, was not implemented and remains for future work.

\subsection{Evaluation against the aims}
These findings address the objectives mentioned in (chapter \ref{chap:introduction}). The first objective was to assess the quality of denoising; the finding is that lower noise does not indicate better spike preservation. The second objective was to determine whether denoising improves detection, and denoising did not reduce the false positives burden when compared to detection on raw signal. The third was to determine whether denoising improves classification, and the results show that classification remained at chance before and after denoising. The fourth was to assess consistency across the five channels, and all findings were consistent across the channels. 
\bigskip

\noindent
Nonetheless, these findings need to be interpreted with caution, because the held-out evaluation set was underpowered, as each channel provided only a small number of evaluation targets due to only using 15 minutes window for analysis; therefore, the at-chance results should be interpreted as consistent with chance rather than conclusively proven results. However, the consistent finding across the five channels strengthens the conclusion. 
\bigskip

\noindent
A notable outcome of this dissertation is that denoising is usually applied as a preprocessing step to improve the quality of the signal, did not improve and, in places, worsened performance for downstream detection and classification. This is most likely because these false positives are actual spikes from other fibres rather than noise, which denoising cannot remove; and also denoising, as seen, distorts and attenuates the spikes, which can reduce the needed information for the downstream stages. So, a lower noise signal does not necessarily mean a better signal for spike sorting. This suggests that denoising should be evaluated by its effect on the later stages of the spike sorting pipeline rather than noise reduction alone. 
\section{Future work}
These findings on single channels are a strong motivation for further and future work. Because the target spikes appear on multiple channels of the recording array, where the spatial signature of each spike is across the channels and its relative amplitude on each electrode, the fibres could be distinguished, unlike single-channel \cite{harris_accuracy_2000, gray_tetrodes_1995}. So, extending this analysis to a multichannel approach and exploiting this spatial information is a natural next step.  
\bigskip

\noindent
Two further directions could be raised from the limitation mentioned above for either single-channel or multi-channel analysis. First, the latency window choice restricted the analysis to a 15-minute window. Over longer recording windows, fibres are more densely active, and activity-dependent slowing causes the latency to shift and spread, so the fixed latency used here cannot capture this spread \cite{serra_activitydependent_1999}. Therefore, a latency tracking approach that follows the moving latency track using track-detection algorithms such as track correlation or activity-dependent shift constraints that follow the ADS-driven latency change over time would allow the method to extend to longer and more active recording periods. This paper by \cite{turnquist_automated_2016} uses automated detection of latency tracks using track correlation, while \cite{troglio_hybrid_2026} uses constraint detection to intervals showing activity-dependent latency shift. Second, which is linked to the first point, a higher-powered evaluation would tighten the confidence intervals and allow the at-chance finding to be tested more conclusively. 
\bigskip

\noindent
Finally, since this work does not perform full spike sorting, but rather binary classification, clustering the detected spikes into multiple distinct units rather than only separating the target from non-target is a natural extension. However, the at-chance results here suggest that full single-channel sorting would have the same limitation, since separating will still require separation by a unit's waveform shape. So, a full sorting might succeed with multi-channel spatial information, which provides a distinguishable signature between the units that a single-channel does not.

Overall, this work shows that denoising on a single channel does not overcome the limitation of the overlap and the waveform similarity between the target fibre and co-active fibres; separating them would need the spatial information.  
% =============================================================================

% Finally, after the main matter, the back matter is specified.  This is
% typically populated with just the bibliography.  LaTeX deals with these
% in one of two ways, namely
%
% - inline, which roughly means the author specifies entries using the 
%   \bibitem macro and typesets them manually, or
% - using BiBTeX, which means entries are contained in a separate file
%   (which is essentially a database) then imported; this is the 
%   approach used below, with the databased being dissertation.bib.
%
% Either way, the each entry has a key (or identifier) which can be used
% in the main matter to cite it, e.g., \cite{X}, \cite[Chapter 2}{Y}.

\backmatter

\bibliography{references.bib}

% -----------------------------------------------------------------------------

% The dissertation concludes with a set of (optional) appendices; these are 
% the same as chapters in a sense, but once signalled as being appendices via
% the associated macro, LaTeX manages them appropriately.

\appendix

\chapter{Appendix}
\label{appx:Appendix}
\begin{table}[htbp]
    \centering
    \begin{tabular}{lllccccccc}
    Classifier & Channel & Densoier & TP & FP & FN & TN & Precision & Recall & F1 \\
    \midrule
    \multirow{15}{*}{One-class}
    & \multirow{3}{*}{9} & Wavelet & 13 & 36 & 3 & 7  & 0.265 & 0.812 & 0.400 \\
    & & Dense & 11 & 30 & 4 & 15 & 0.268 & 0.733 & 0.393 \\
    & & Conv & 8 & 33 & 9 & 13 & 0.195 & 0.471 & 0.276 \\
    \cmidrule{2-10}
    & \multirow{3}{*}{14} & Wavelet & 14 & 37 & 2 & 6  & 0.275 & 0.875 & 0.418 \\
     & & Dense & 8 & 29 & 2 & 16 & 0.216 & 0.800 & 0.340 \\
     & & Conv & 8 & 28 & 7 & 22 & 0.222 & 0.533 & 0.314 \\
    \cmidrule{2-10}
     & \multirow{3}{*}{19} & Wavelet & 11 & 40 & 3 & 4  & 0.216 & 0.786 & 0.338 \\
     & & Dense & 6 & 40 & 2 & 16 & 0.130 & 0.750 & 0.222 \\
     & & Conv & 8 & 33 & 9 & 13 & 0.195 & 0.471 & 0.276 \\
     \cmidrule{2-10}
     & \multirow{3}{*}{25} & Wavelet & 11 & 31 & 1 & 9  & 0.262 & 0.917 & 0.407 \\
     & & Dense & 20 & 33 & 1 & 7 & 0.377 & 0.952 & 0.541 \\
     & & Conv & 12 & 38 & 0 & 11 & 0.240 & 1.000 & 0.387 \\
    \cmidrule{2-10}
     & \multirow{3}{*}{30} & Wavelet & 6  & 37 & 1 & 7  & 0.140 & 0.857 & 0.240 \\
     & & Dense & 10 & 35 & 4 & 16 & 0.222 & 0.714 & 0.339 \\
     & & Conv & 9 & 28 & 2 & 23 & 0.243 & 0.818 & 0.375 \\
    \midrule
    \multirow{15}{*}{Two-class}
     & \multirow{3}{*}{9} & Wavelet & 9  & 25 & 7 & 18 & 0.265 & 0.562 & 0.360 \\
     & & Dense & 8  & 20 & 7 & 25 & 0.286 & 0.533 & 0.372 \\
     & & Conv & 10 & 24 & 6 & 25 & 0.294 & 0.625 & 0.400 \\
    \cmidrule{2-10}
     & \multirow{3}{*}{14} & Wavelet & 7  & 22 & 9 & 21 & 0.241 & 0.438 & 0.311 \\
     & & Dense   & 7 & 26 & 3 & 19 & 0.212 & 0.700 & 0.326 \\
     & & Conv & 8 & 27 & 7 & 23 & 0.229 & 0.533 & 0.320 \\
    \cmidrule{2-10}
     & \multirow{3}{*}{19} & Wavelet & 8  & 22 & 6  & 22 & 0.267 & 0.571 & 0.364 \\
     & & Dense & 4 & 27 & 4 & 29 & 0.129 & 0.500 & 0.205 \\
     & & Conv & 5 & 20 & 12 & 26 & 0.200 & 0.294 & 0.238 \\
    \cmidrule{2-10}
     & \multirow{3}{*}{25} & Wavelet & 5 & 24 & 7 & 16 & 0.172 & 0.417 & 0.244 \\
     & & Dense & 7 & 24 & 14 & 16 & 0.226 & 0.333 & 0.269 \\
     & & Conv & 7 & 26 & 5 & 23 & 0.212 & 0.583 & 0.311 \\
    \cmidrule{2-10}
     & \multirow{3}{*}{30} & Wavelet & 2  & 24 & 5  & 20 & 0.077 & 0.286 & 0.121 \\
     & & Dense & 4 & 30 & 10 & 21 & 0.118 & 0.286 & 0.167 \\
     & & Conv & 10 & 32 & 1 & 19 & 0.238 & 0.909 & 0.377 \\
    \bottomrule
    \end{tabular}
    \caption{Full classification metrics after denoising per channel, denoiser and classifier on the held-out evaluation. The TP, FP, FN, and TN are the confusion matrix counts; the precision, recall and F1 are derived from them.}
    \label{tab:confusion matrix classification after denoising}
\end{table}



% =============================================================================

\end{document}
